\section{引言}

场空间中的流是一种有趣的工具，它可能使我们对于高度非线性量子场论（如 QCD）所
描述的物理机制获得新的认识。本文研究的 $\SUthree$ 规范场的流 $B_{\mu}(t,x)$ 由
如下方程定义
\begin{align}
  \dot{B}_{\mu} &= D_{\nu}G_{\nu\mu},
  &
  \left.B_{\mu}\right|_{t=0} &= A_{\mu},
  \label{eq:floweq}\\
  G_{\mu\nu} &= \partial_{\mu}B_{\nu}-\partial_{\nu}B_{\mu}
  +[B_{\mu},B_{\nu}],
  &
  D_{\mu} &= \partial_{\mu}+[B_{\mu},\,\cdot\,],
\end{align}
其中 $A_{\mu}$ 是 QCD 中的基本规范场
（未加解释的记号见附录~\ref{app:A}；对“流时间”$t$ 的导数用一点表示）。显然，随
着 $t$ 的增大，只要不出现奇异性，流方程~\eqref{eq:floweq} 便沿最速下降方向驱动
规范场趋向 Yang--Mills 作用量的驻点。

在格点 QCD 中，规范场 $U(x,\mu)$ 的作用量最简单的选择是 Wilson
作用量~\cite{Wilson}
\begin{equation}
  \Sw(U)=\frac{1}{g_0^2}\sum_p
  \Re\tr\{1-U(p)\},
  \label{eq:wact}
\end{equation}
其中 $g_0$ 是裸耦合，$p$ 遍历格点上所有有向方格（plaquette），$U(p)$ 表示绕 $p$
一周链变量的乘积。与之相应的格点规范场的流 $V_t(x,\mu)$（即“威尔逊流”）由方程
\begin{equation}
  \dot{V}_t(x,\mu)=-g_0^2\left\{\partial_{x,\mu}\Sw(V_t)\right\}V_t(x,\mu),
  \qquad
  \left.V_t(x,\mu)\right|_{t=0}=U(x,\mu),
  \label{eq:wflow}
\end{equation}
定义，式中 $\partial_{x,\mu}$ 表示对链变量 $V_t(x,\mu)$ 的自然 $\suthree$ 型微分
算符（见附录~\ref{app:A}）。在有穷格点上，威尔逊流在一切正、负时间 $t$ 处的存在
性、唯一性与光滑性都得到严格保证~\cite{ArnoldI}。此外，由式~\eqref{eq:wflow}
立即可知，作用量 $\Sw(V_t)$ 是 $t$ 的单调下降函数。因此流趋于对场起光滑化作用；
事实上，它由无穷小的“stout 链涂抹”（stout link smearing）步骤生成~\cite{Stout}。

威尔逊流此前在文献~\cite{TrivMaps} 中于场空间平凡化映射的背景下出现过。本文并不
要求读者熟悉该文，但那里得到的若干数学结果将在此再次使用。下文的一个重要目标，
是考察由正流时间处的规范场构造的局域观测量，其期望值是否可以预期具有良定义的连
续极限。文中通过直接在连续理论中利用维数正规化完成的一个一圈阶微扰计算示例，以及
对三个晶格间距值的 $\SUthree$ 规范理论的数值研究，为该极限的存在提供了证据。随后
讨论了威尔逊流的两个应用：一是格点 QCD 中的定标（scale-setting）问题，二是当晶格
间距趋于零时拓扑（瞬子）扇区究竟如何涌现的问题。
